% Luc Destombe --- licence CC BY-NC-SA 4.0
% https://creativecommons.org/licenses/by-nc-sa/4.0/deed.fr
% Fichier autonome : compiler avec pdflatex, deux fois (signets, nombre de pages).
\def\lyceeversion{1.0}
\documentclass[a4paper,10pt]{article}
\makeatletter
\newif\iflycee@niveaux
\lycee@niveauxtrue
\newif\iflycee@cm
\newif\iflycee@palatino

\RequirePackage[utf8]{inputenc}
\RequirePackage[T1]{fontenc}
\RequirePackage[french]{babel}
\RequirePackage{etoolbox}

\newcommand{\ShorthandsOff}{\@ifundefined{shorthandoff}{}{\shorthandoff{;:!?}}}
\newcommand{\ShorthandsOn}{\@ifundefined{shorthandon}{}{\shorthandon{;:!?}}}
\ShorthandsOff
\AtBeginDocument{\ShorthandsOn}
\AtBeginEnvironment{tikzpicture}{\ShorthandsOff}

\RequirePackage{geometry}
\RequirePackage{fancyhdr}
\RequirePackage{lastpage}
\RequirePackage{multicol}
\RequirePackage{array}
\RequirePackage{enumitem}
\RequirePackage{xcolor}
\RequirePackage{needspace}

\RequirePackage{amsmath,amssymb}
\iflycee@cm\else\RequirePackage{fourier}\fi
\DeclareMathAlphabet{\mathcal}{OMS}{cmsy}{m}{n}
\renewcommand{\ttdefault}{lmtt}
\begingroup\catcode`\_=\active \catcode`\^=\active
\global\let_\sb \global\let^\sp
\endgroup
\AtBeginDocument{\catcode`\_=12 \mathcode`\_="8000
  \catcode`\^=12 \mathcode`\^="8000 }
\newlength{\retraitcalculs}
\setlength{\retraitcalculs}{2em}

\RequirePackage{tikz}
\usetikzlibrary{arrows.meta,calc,babel,shapes.misc}
\RequirePackage[most]{tcolorbox}
\tcbuselibrary{skins,breakable}

\definecolor{coulDef}{HTML}{1F4E79}
\definecolor{coulProp}{HTML}{7030A0}
\definecolor{coulRem}{HTML}{C55A11}
\definecolor{coulEx}{HTML}{2E7D32}
\definecolor{coulExo}{HTML}{B03A2E}
\definecolor{coulPart}{HTML}{1F4E79}
\definecolor{coulMeth}{HTML}{1B6E4F}
\definecolor{coulCourbe}{HTML}{0B5ED7}
\definecolor{coulCourbeB}{HTML}{C00000}
\definecolor{coulCourbeC}{HTML}{2E7D32}
\definecolor{coulGrille}{HTML}{8C8C8C}
\definecolor{coulRep}{HTML}{1B6E4F}
\definecolor{coulBar}{HTML}{2E7D32}

\newcommand{\R}{\mathbb{R}}
\newcommand{\Rs}{\mathbb{R}^{*}}
\newcommand{\N}{\mathbb{N}}
\newcommand{\Ns}{\mathbb{N}^{*}}
\newcommand{\Z}{\mathbb{Z}}
\newcommand{\Q}{\mathbb{Q}}
\renewcommand{\le}{\leqslant}
\renewcommand{\ge}{\geqslant}
\renewcommand{\approx}{\simeq}
\newcommand{\vect}[1]{\overrightarrow{#1}}
\newcommand{\norme}[1]{\left\lVert #1 \right\rVert}
\newcommand{\intervalle}[2]{\left[#1\,;#2\right]}
\RequirePackage{cancel}
\renewcommand{\CancelColor}{\color{coulExo}}
\newcommand{\fois}[1]{\times{\color{coulExo}#1}}
\newcommand{\barre}[1]{\cancel{#1}}
\newlength{\lycee@eqh}\newlength{\lycee@eqd}
\newsavebox{\lycee@eqa}\newsavebox{\lycee@eqb}\newsavebox{\lycee@eqc}
\newcommand{\lycee@eqcol}[1]{\hbox{$\renewcommand{\arraystretch}{1.35}%
  \begin{array}[t]{@{}>{\displaystyle}r@{}>{\displaystyle{}}l@{}}#1\end{array}$}}
\newcommand{\lycee@eqmesure}[1]{%
  \ifdim\ht#1>\lycee@eqh\setlength{\lycee@eqh}{\ht#1}\fi
  \ifdim\dp#1>\lycee@eqd\setlength{\lycee@eqd}{\dp#1}\fi}
\newcommand{\lycee@eqtrait}{\hspace{1.2em}ou\hspace{1.2em}}
\newcommand{\lycee@eqiff}{$\iff$\hspace{0.8em}}
\newcommand{\eqproduit}[2]{\par\smallskip\noindent C'est une équation produit :
  $A\times B=0 \iff A=0$ ou $B=0$.\par\nopagebreak\smallskip
  \begingroup\setlength{\lycee@eqh}{0pt}\setlength{\lycee@eqd}{0pt}%
  \setbox\lycee@eqa=\lycee@eqcol{#1}\setbox\lycee@eqb=\lycee@eqcol{#2}%
  \lycee@eqmesure\lycee@eqa\lycee@eqmesure\lycee@eqb%
  \noindent\hspace*{\retraitcalculs}\lycee@eqiff\usebox\lycee@eqa\lycee@eqtrait\usebox\lycee@eqb\endgroup\par\smallskip}
\newcommand{\eqproduittrois}[3]{\par\smallskip\noindent C'est une équation produit :
  $A\times B\times C=0 \iff A=0$ ou $B=0$ ou $C=0$.\par\nopagebreak\smallskip
  \begingroup\setlength{\lycee@eqh}{0pt}\setlength{\lycee@eqd}{0pt}%
  \setbox\lycee@eqa=\lycee@eqcol{#1}\setbox\lycee@eqb=\lycee@eqcol{#2}\setbox\lycee@eqc=\lycee@eqcol{#3}%
  \lycee@eqmesure\lycee@eqa\lycee@eqmesure\lycee@eqb\lycee@eqmesure\lycee@eqc%
  \noindent\hspace*{\retraitcalculs}\lycee@eqiff\usebox\lycee@eqa\lycee@eqtrait\usebox\lycee@eqb\lycee@eqtrait\usebox\lycee@eqc\endgroup
  \par\smallskip}

\tikzset{grille/.style={coulGrille,line width=0.4pt}}
\newcommand{\quadrillage}[4]{\draw[grille,step=1] (#1,#2) grid (#3,#4);}
\tikzset{
  croix/.style={cross out,draw,line width=0.6pt,inner sep=0pt,minimum size=4pt},
  croixdroite/.style={croix,rotate=45,line width=0.8pt},
}
\newcommand{\pt}[3]{\path (#1) node[croix]{} node[#2,font=\small,inner sep=2.5pt]{$#3$};}

\newcommand{\lycee@repere}[4]{%
  \draw[grille,step=1] (#1,#2) grid (#3,#4);
  \draw[-{Stealth[length=2mm]},thick] (#1,0)--({#3+0.4},0) node[below left=-1pt]{$x$};
  \draw[-{Stealth[length=2mm]},thick] (0,#2)--(0,{#4+0.4}) node[below left=-1pt]{$y$};
  \node[below left,font=\scriptsize,inner sep=1.5pt] at (0,0) {$0$};}
\newcommand{\lycee@gradx}[1]{\foreach \k/\l in {#1}{%
  \draw (\k,0.07)--(\k,-0.07) node[below,font=\scriptsize,inner sep=1.5pt]{$\l$};}}
\newcommand{\lycee@grady}[1]{\foreach \k/\l in {#1}{%
  \draw (0.07,\k)--(-0.07,\k) node[left,font=\scriptsize,inner sep=1.5pt]{$\l$};}}
\newcommand{\lycee@intff}[2]{\left[#1\,;#2\right]}
\newcommand{\lycee@intfo}[2]{\left[#1\,;#2\right[}
\newcommand{\lycee@intof}[2]{\left]#1\,;#2\right]}
\newcommand{\lycee@intoo}[2]{\left]#1\,;#2\right[}
\newcommand{\lycee@ens}[1]{\left\{#1\right\}}
\AtBeginDocument{%
  \providecommand{\repere}{\lycee@repere}%
  \providecommand{\gradx}{\lycee@gradx}%
  \providecommand{\grady}{\lycee@grady}%
  \providecommand{\Cf}{\mathcal{C}\sb{f}}%
  \providecommand{\Cg}{\mathcal{C}\sb{g}}%
  \providecommand{\intff}{\lycee@intff}%
  \providecommand{\intfo}{\lycee@intfo}%
  \providecommand{\intof}{\lycee@intof}%
  \providecommand{\intoo}{\lycee@intoo}%
  \providecommand{\ens}{\lycee@ens}}

\newlist{questions}{enumerate}{3}
\setlist[questions,1]{label=\arabic*\textdegree),leftmargin=*,itemsep=2pt,topsep=3pt}
\setlist[questions,2]{label=\alph*),leftmargin=*,itemsep=1pt,topsep=2pt}
\setlist[questions,3]{label=\roman*),leftmargin=*,itemsep=1pt,topsep=1pt}
\setlist[itemize]{label=\textbullet,leftmargin=*,itemsep=2pt,topsep=3pt}

\newcommand{\besoinplace}[1]{\par\penalty\@M\begingroup
  \dimen@=#1\relax\dimen@ii=\pagegoal\advance\dimen@ii-\pagetotal
  \ifdim\dimen@>\dimen@ii\ifdim\dimen@ii>\z@\vfil\fi\break\fi\endgroup}

\newcommand{\lycee@pied}{\thepage/\pageref{LastPage}}
\newcommand{\entete}[2]{%
  \pagestyle{fancy}\fancyhf{}%
  \fancyhead[L]{\footnotesize\color{coulPart}\textbf{#1}}%
  \fancyhead[R]{\footnotesize\color{coulPart}\textbf{#2}}%
  \fancyfoot[C]{\footnotesize\lycee@pied}%
  \renewcommand{\headrulewidth}{0.4pt}%
  \renewcommand{\headrule}{\hbox to\headwidth{%
    \color{coulPart!40}\leaders\hrule height \headrulewidth\hfill}}}

\RequirePackage{ccicons}
\newcommand{\lycee@licence}{{\scriptsize\color{gray}%
  \copyright\ Luc Destombe --- \ccbyncsa\ CC BY-NC-SA 4.0}}
\newif\iflycee@licence \lycee@licencetrue
\newcommand{\sanslicence}{\lycee@licencefalse}
\AtBeginDocument{\iflycee@licence\AddToHookNext{shipout/foreground}{%
  \put(0,0){\rlap{\kern\dimexpr1in+\hoffset+\oddsidemargin+\textwidth\relax
    \llap{\raisebox{-\dimexpr1in+\voffset+\topmargin+\headheight+\headsep
      +\textheight+\footskip\relax}{\lycee@licence}}}}}\fi}

\newcommand{\lycee@site}{faire-des-maths.fr}
\newcommand{\lycee@version}{{\scriptsize\color{gray}\lycee@site
  \ifdefined\lyceeversion\ --- Version~\lyceeversion\fi}}
\AtBeginDocument{\iflycee@licence\AddToHookNext{shipout/foreground}{%
  \put(0,0){\rlap{\kern\dimexpr1in+\hoffset+\oddsidemargin\relax
    \raisebox{-\dimexpr1in+\voffset+\topmargin+\headheight+\headsep
      +\textheight+\footskip\relax}{\lycee@version}}}}\fi}

\newcommand{\formulesserrees}{%
  \setlength{\abovedisplayskip}{3pt}\setlength{\belowdisplayskip}{3pt}%
  \setlength{\abovedisplayshortskip}{2pt}\setlength{\belowdisplayshortskip}{3pt}}

\setlength{\parindent}{0pt}

\RequirePackage[implicit=false,hidelinks,pdfencoding=auto,psdextra,bookmarksopen,
  bookmarksopenlevel=1,pdfcreator={},pdfproducer={}]{hyperref}
\RequirePackage{bookmark}
\hypersetup{pdfauthor={Luc Destombe},pdfkeywords={CC BY-NC-SA 4.0}}
\newcounter{lycee@signet}
\newcommand{\signet}[3][]{\stepcounter{lycee@signet}%
  \edef\lycee@dest{\if\relax\detokenize{#1}\relax
    signet.\arabic{lycee@signet}\else#1\fi}%
  \hypertarget{\lycee@dest}{}\bookmark[level=#2,dest=\lycee@dest]{#3}}
\pdfstringdefDisableCommands{%
  \def\R{R}\def\N{N}\def\Z{Z}\def\Q{Q}\def\vect#1{#1}\def\Cf{Cf}\def\Cg{Cg}\def\fois#1{×#1}%
  \def\barre#1{#1}\def\intervalle#1#2{[#1;#2]}\def\intff#1#2{[#1;#2]}%
  \def\textsuperscript#1{#1}\def\dfrac#1#2{#1/#2}\def\frac#1#2{#1/#2}%
  \def\vec#1{#1⃗}}%
\RequirePackage[official]{eurosym}

\geometry{top=0.8cm,bottom=1.3cm,left=1cm,right=1cm,footskip=0.6cm}

\pagestyle{fancy}
\fancyhf{}
\renewcommand{\headrulewidth}{0pt}
\fancyfoot[C]{\footnotesize Page \thepage{} / \pageref{LastPage}}

\newcommand{\titrefiche}[2]{%
  \begin{center}
    \Large\textbf{#1}\\[2pt]
    \large\textbf{#2}\\[2pt]
  \end{center}
  \vspace{0.10cm}}

\setlength{\columnseprule}{0.4pt}
\setlength{\columnsep}{15pt}
\renewcommand{\columnseprulecolor}{\color{black!45}}
\setlength{\multicolsep}{2pt}

\newcounter{partiefiche}
\renewcommand{\thepartiefiche}{\Roman{partiefiche}}
\newif\ifapresbandeau
\newcommand{\lycee@bandeau}[1]{%
  \par\penalty-600
  \vspace{0.08cm}%
  \noindent\colorbox{coulPart}{%
    \parbox{\dimexpr\linewidth-2\fboxsep\relax}{%
      \color{white}\bfseries\raggedright\strut#1\strut}}%
  \nopagebreak\par\nopagebreak
  \vspace{0.06cm}\nopagebreak
  \apresbandeautrue}
\newcommand{\partiefiche}[1]{%
  \refstepcounter{partiefiche}%
  \lycee@bandeau{\signet{1}{\thepartiefiche{} --- #1}\thepartiefiche{} --- #1}}
\newcommand{\partiefichesansnum}[1]{\lycee@bandeau{\signet{1}{#1}#1}}

\newcounter{exo}
\newlength{\espaceexo}\setlength{\espaceexo}{7pt}
\newlength{\espacetitre}\setlength{\espacetitre}{2pt}
\newcommand{\exo}[1][\relax]{%
  \par
  \ifapresbandeau\apresbandeaufalse\else\penalty-150\addvspace{\espaceexo}\fi
  \refstepcounter{exo}%
  \noindent\signet[exercice-\arabic{exo}]{2}{Exercice \arabic{exo}}%
  \textbf{\color{coulExo}\underline{Exercice \theexo}}%
  \ifx\relax#1\relax\else\textbf{ : #1}\fi
  \par\nopagebreak\vspace{\espacetitre}\nopagebreak
  \ignorespaces}

\setlist[enumerate]{nosep,leftmargin=*,itemsep=2pt,topsep=2pt}
\setlist[itemize]{nosep,leftmargin=*,topsep=2pt}
\newcommand{\choix}[3]{\par\hspace*{1em}%
  \makebox[0.3\linewidth][l]{a) #1}\makebox[0.3\linewidth][l]{b) #2}c) #3}
\newcommand{\deux}[2]{\par\hspace*{0.6em}%
  \makebox[0.48\linewidth][l]{#1}#2}
\newcommand{\trois}[3]{\par\hspace*{0.6em}%
  \makebox[0.32\linewidth][l]{#1}\makebox[0.32\linewidth][l]{#2}#3}

  \renewcommand{\exo}[1][\relax]{%
    \ifapresbandeau\apresbandeaufalse\else\par\penalty-150\fi
    \refstepcounter{exo}%
    \vspace{0.05cm}%
    \noindent\signet[exercice-\arabic{exo}]{2}{Exercice \arabic{exo}}%
    \textbf{\color{coulExo}\underline{Exercice \theexo}}%
    \ifx\relax#1\relax\else\textbf{ : #1}\fi%
    \par\nopagebreak
    \ignorespaces}
  \newcommand{\rappel}[1]{%
    \par\vspace{0.06cm}\nopagebreak
    \noindent\fcolorbox{black!35}{black!5}{%
      \parbox{\dimexpr\linewidth-2\fboxsep-2\fboxrule\relax}{%
        \small\textbf{Rappel.} #1}}%
    \par\vspace{0.06cm}\nopagebreak}
  \setlist[enumerate]{nosep,leftmargin=*}
  \setlist[itemize]{nosep,leftmargin=*}
  \setlength{\multicolsep}{3pt plus 1pt minus 1pt}
  \setlength{\premulticols}{10pt}
  \setlength{\postmulticols}{10pt}
  \newlist{expr}{enumerate}{1}
  \setlist[expr]{label={\Alph*\ $=$},nosep,leftmargin=*,itemsep=0pt}

\renewenvironment{center}{\par\addvspace{3pt}\centering}{\par\addvspace{3pt}}
\formulesserrees
\AtBeginDocument{\formulesserrees}
\clubpenalty=10000
\widowpenalty=10000
\displaywidowpenalty=10000
\brokenpenalty=10000
\setlength{\parskip}{0pt}

\RequirePackage{ragged2e}
\setlength{\RaggedRightRightskip}{0pt plus 1fil}
\AtBeginDocument{\RaggedRight\binoppenalty=10000 \relpenalty=10000 }
\g@addto@macro\@parboxrestore{\RaggedRight}
\makeatother
\geometry{top=0.8cm,bottom=1.3cm,left=1cm,right=1cm,footskip=0.7cm,headheight=12pt}

\newtcblisting{python}{%
  listing only,enhanced,breakable,
  colback=black!3,colframe=black!35,boxrule=0.5pt,arc=1pt,
  left=6pt,right=6pt,top=2pt,bottom=2pt,
  listing options={language=Python,basicstyle=\ttfamily\small,
    keywordstyle=\color{coulMeth}\bfseries,columns=fullflexible,
    keepspaces=true,showstringspaces=false,aboveskip=0pt,belowskip=0pt}}

\begin{document}
\titrefiche{Terminale spécialité --- Loi binomiale}{Fiche d'exercices du chapitre}

\begin{multicols}{2}

\partiefiche{Rappels : probabilités conditionnelles}

\exo[Compléter un arbre, sans calculatrice]
Soit $A$ et $B$ deux événements tels que $P(A)=0{,}3$, $P_{A}(B)=0{,}6$ et
$P_{\overline{A}}(B)=0{,}2$.
\begin{enumerate}
  \item Représenter la situation par un arbre pondéré complet.
  \item Calculer $P(A\cap B)$, puis $P(B)$.
  \item En déduire $P_{B}(A)$.
  \item Les événements $A$ et $B$ sont-ils indépendants ?
\end{enumerate}

\exo[Les terrains d'une salle de badminton]
Une salle de badminton classe ses heures d'ouverture en heures creuses et heures
pleines ; $65\,\%$ des heures sont creuses. Pendant une heure creuse, $30\,\%$ des
terrains sont occupés ; pendant une heure pleine, $80\,\%$ le sont. On choisit au hasard
un terrain et une heure d'ouverture, et on note $C$ : \og l'heure est creuse\fg{} et $T$ :
\og le terrain est occupé\fg.
\begin{enumerate}
  \item Traduire les trois pourcentages de l'énoncé par des probabilités.
  \item Représenter la situation par un arbre pondéré.
  \item Calculer la probabilité que l'heure soit creuse et le terrain occupé.
  \item Calculer la probabilité que le terrain soit occupé.
  \item Le terrain est occupé. Calculer la probabilité que l'heure soit pleine, arrondie
        au millième.
\end{enumerate}

\exo[Deux capteurs indépendants]
Un drone d'observation utilise deux capteurs d'altitude, qui tombent en panne
indépendamment l'un de l'autre. On note $A$ : \og le capteur A est en panne\fg, de
probabilité $0{,}01$, et $B$ : \og le capteur B est en panne\fg, de probabilité $0{,}03$.
\begin{enumerate}
  \item Calculer la probabilité que les deux capteurs soient en panne.
  \item Calculer la probabilité qu'au moins un des deux capteurs soit en panne.
  \item Le capteur A est en panne. Quelle est la probabilité que le capteur B le soit
        aussi ?
\end{enumerate}

\exo[D'après bac 2025 et 2026 : vrai ou faux ?]
Pour chaque affirmation, indiquer si elle est vraie ou fausse, en justifiant.
\begin{enumerate}
  \item $52\,\%$ des musiciens professionnels jouent de la Pop ; parmi eux, $32\,\%$ sont
        des femmes ; $20\,\%$ des musiciens professionnels sont des femmes. On choisit
        au hasard un musicien professionnel.

        \emph{Affirmation} : si la personne choisie est une femme, la probabilité
        qu'elle joue de la Pop est de $0{,}832$.
  \item Dans un musée, la probabilité qu'un visiteur achète son billet en ligne est
        $0{,}7$ ; la probabilité qu'il choisisse l'audioguide sachant qu'il a acheté son
        billet en ligne est $0{,}8$ ; la probabilité qu'il visite sans audioguide est
        $0{,}32$.

        \emph{Affirmation} : la probabilité qu'un visiteur ne prenne pas l'audioguide
        sachant qu'il a acheté son billet au guichet est supérieure à deux tiers.
  \item Soit $A$ et $B$ deux événements tels que $P(A)=\frac{2}{5}$,
        $P_{A}(B)=\frac{1}{4}$ et $P_{\overline{A}}(B)=\frac{3}{10}$.

        \emph{Affirmation} : $P_{B}(\overline{A})=\dfrac{9}{50}$.
\end{enumerate}

\partiefiche{Rappels : variable aléatoire, loi, espérance}

\exo[Une roue de loterie]
À une kermesse, on mise $4$\,€ et on fait tourner une roue partagée en $8$ secteurs
égaux : un secteur fait gagner $20$\,€, trois secteurs font gagner $5$\,€, les quatre
autres ne font rien gagner. On note $X$ le gain algébrique du joueur.
\begin{enumerate}
  \item Déterminer la loi de probabilité de $X$.
  \item Calculer $E(X)$. Le jeu est-il favorable au joueur ?
  \item Calculer $V(X)$, puis $\sigma(X)$ arrondi au centième.
  \item Les organisateurs veulent un jeu équitable. Quelle mise doivent-ils demander ?
\end{enumerate}

\exo[La recette d'une salle de badminton]
Une salle de badminton loue ses terrains $10$\,€ l'heure pleine et $6$\,€ l'heure creuse.
On note $X$ la recette, en euros, d'un terrain pendant une heure choisie au hasard. Une
étude donne la loi de $X$ :
\begin{center}
\renewcommand{\arraystretch}{1.4}
\begin{tabular}{|c|c|c|c|}
\hline
$x_{i}$ & $0$ & $6$ & $10$\\\hline
$P(X=x_{i})$ & $0{,}525$ & $0{,}195$ & $0{,}28$\\\hline
\end{tabular}
\end{center}
\begin{enumerate}
  \item Calculer $E(X)$ et interpréter le résultat.
  \item La salle compte $10$ terrains et ouvre $70$ heures par semaine. Estimer sa
        recette hebdomadaire moyenne.
  \item Calculer $\sigma(X)$, arrondi au centième.
  \item Une autre salle a une recette horaire d'espérance $3$\,€ et d'écart-type $4$\,€.
        Laquelle rapporte le plus en moyenne ? Laquelle a les recettes les plus
        régulières ?
\end{enumerate}

\exo[D'après bac 2025, Métropole : vrai ou faux ?]
Dans un musée, un visiteur muni d'un audioguide choisit l'un de trois parcours, de
durées $50$~min, $1$~h~$20$~min et $1$~h~$40$~min. La durée $X$ du parcours suit la loi
ci-dessous.
\begin{center}
\renewcommand{\arraystretch}{1.4}
\begin{tabular}{|c|c|c|c|}
\hline
$x_{i}$ & $50$~min & $1$~h~$20$~min & $1$~h~$40$~min\\\hline
$P(X=x_{i})$ & $0{,}1$ & $0{,}6$ & $0{,}3$\\\hline
\end{tabular}
\end{center}
\emph{Affirmation} : l'espérance de $X$ est $77$ minutes. Vraie ou fausse ? Justifier.

\partiefiche{Succession d'épreuves indépendantes}

\exo[Un dé, puis une pièce]
On lance un dé équilibré à quatre faces numérotées de $1$ à $4$, puis une pièce
équilibrée dont les faces portent les numéros $1$ et $2$. Une issue est un couple, comme
$(3\,;2)$.
\begin{enumerate}
  \item Combien cette expérience a-t-elle d'issues ? Représenter la situation par un
        arbre pondéré.
  \item Déterminer la loi de probabilité sur l'ensemble des issues.
  \item Calculer la probabilité d'obtenir deux numéros pairs.
  \item Calculer la probabilité que la somme des deux numéros soit égale à $4$.
\end{enumerate}

\exo[Trois urnes]
On tire une boule dans chacune de trois urnes, l'une après l'autre :
\begin{itemize}
  \item l'urne 1 contient $2$ boules rouges et $3$ boules bleues ;
  \item l'urne 2 contient $1$ boule rouge et $3$ boules vertes ;
  \item l'urne 3 contient $1$ boule rouge et $4$ boules jaunes.
\end{itemize}
On note $R_{1}$, $R_{2}$, $R_{3}$ les événements \og la boule tirée dans l'urne 1 (2, 3)
est rouge\fg. On gagne si l'on obtient au moins deux boules rouges.
\begin{enumerate}
  \item Représenter la situation par un arbre pondéré.
  \item Calculer la probabilité d'obtenir trois boules rouges.
  \item Calculer la probabilité de gagner.
  \item Calculer la probabilité d'obtenir au moins une boule rouge.
\end{enumerate}

\exo[Le contrôle d'un vélo électrique]
À la sortie d'une usine, un vélo électrique peut présenter un défaut de batterie (probabilité
$0{,}03$), de moteur ($0{,}05$) ou de freinage ($0{,}1$). Ces défauts apparaissent
indépendamment les uns des autres. On note $B$, $M$ et $F$ ces trois événements, et on
choisit un vélo au hasard.
\begin{enumerate}
  \item Représenter la situation par un arbre pondéré à trois niveaux.
  \item Calculer la probabilité que le vélo présente les trois défauts.
  \item Calculer $P\big((\overline{B}\,;M\,;\overline{F})\big)$ et traduire cette issue par
        une phrase.
  \item Calculer la probabilité que le vélo ne présente aucun défaut, puis la probabilité
        qu'il en présente au moins un. Arrondir au millième.
\end{enumerate}

\exo[Répondre au hasard]
Un questionnaire comporte cinq questions ; pour chacune, deux réponses sont proposées,
l'une juste, l'autre fausse. Un élève répond au hasard à chaque question, ses réponses étant
indépendantes.
\begin{enumerate}
  \item Combien y a-t-il de façons de remplir ce questionnaire ? Que peut-on dire de la
        loi de probabilité sur l'ensemble de ces réponses ?
  \item Calculer, sous forme de fraction irréductible, la probabilité de chaque
        événement :
  \begin{enumerate}[label=\alph*)]
    \item $A$ : \og l'élève donne cinq réponses justes\fg{} ;
    \item $B$ : \og l'élève donne une seule réponse juste\fg{} ;
    \item $C$ : \og l'élève donne au moins deux réponses justes\fg.
  \end{enumerate}
\end{enumerate}

\partiefiche{Épreuve, loi et schéma de Bernoulli}

\exo[Avec ou sans remise]
Une urne contient $4$ boules vertes et $2$ boules rouges. On tire deux boules l'une après
l'autre, et on note $V_{1}$ : \og la première boule est verte\fg{} et $V_{2}$ : \og la
seconde boule est verte\fg.
\begin{enumerate}
  \item Les tirages se font \textbf{sans remise}.
  \begin{enumerate}[label=\alph*)]
    \item Représenter la situation par un arbre pondéré.
    \item Calculer $P(V_{1}\cap V_{2})$, puis $P(V_{2})$.
    \item Les événements $V_{1}$ et $V_{2}$ sont-ils indépendants ?
  \end{enumerate}
  \item Mêmes questions lorsque les tirages se font \textbf{avec remise}.
  \item Dans lequel des deux cas la répétition du tirage d'une boule forme-t-elle un
        schéma de Bernoulli ?
\end{enumerate}

\exo[Schéma de Bernoulli ou non ?]
Pour chaque situation, dire si elle peut être modélisée par un schéma de Bernoulli et,
si oui, en donner les paramètres $n$ et $p$.
\begin{enumerate}[label=\alph*)]
  \item On lance $10$ fois un dé équilibré à six faces ; le succès est \og obtenir un
        $6$\fg.
  \item On tire successivement et sans remise $3$ cartes d'un jeu de $32$ cartes ; le
        succès est \og tirer un as\fg.
  \item On tire $4$ fois, avec remise, une boule dans une urne contenant $3$ boules
        rouges et $5$ boules noires ; le succès est \og la boule est rouge\fg.
  \item On choisit au hasard $20$ personnes dans une grande ville, en assimilant ce choix
        à un tirage avec remise ; $45\,\%$ des habitants ont un abonnement aux transports
        en commun, et le succès est \og la personne a un abonnement\fg.
\end{enumerate}

\exo[Des briques de construction]
Un bac contient $36$ briques de construction jaunes et $64$ vertes. Lina prend une brique
au hasard, regarde sa couleur et la remet dans le bac ; elle recommence trois fois en
tout.
\begin{enumerate}
  \item Décrire l'épreuve de Bernoulli associée à un tirage, en prenant pour succès
        \og la brique est jaune\fg.
  \item Justifier que la situation est un schéma de Bernoulli et donner ses paramètres.
  \item Représenter ce schéma par un arbre pondéré.
  \item Calculer la probabilité que Lina obtienne exactement deux briques jaunes.
  \item Calculer la probabilité qu'elle n'obtienne aucune brique jaune.
\end{enumerate}

\exo[La variance d'une loi de Bernoulli]
Soit $p$ un nombre réel de $\intff{0}{1}$ et $X$ une variable aléatoire qui suit la loi
de Bernoulli de paramètre $p$.
\begin{enumerate}
  \item Pour $p=0{,}2$, calculer $E(X)$, $V(X)$ et $\sigma(X)$.
  \item On rappelle que $V(X)=p(1-p)$. Pour quelle valeur de $p$ la variance de $X$
        est-elle la plus grande ? Quelle est alors sa valeur ?
  \item Pour quelles valeurs de $p$ la variance est-elle nulle ? Interpréter.
\end{enumerate}

\partiefiche{Loi binomiale}

\exo[Calculer avec la formule]
Soit $X$ une variable aléatoire qui suit la loi binomiale $\mathcal{B}(5\,;0{,}6)$.
\begin{enumerate}
  \item Calculer $P(X=0)$, $P(X=2)$ et $P(X=5)$, en écrivant d'abord la formule.
  \item Compléter le tableau de la loi de $X$ (valeurs exactes), puis vérifier que la
        somme des probabilités vaut $1$.
  \begin{center}
  \renewcommand{\arraystretch}{1.4}
  \begin{tabular}{|c|c|c|c|c|c|c|}
  \hline
  $k$ & $0$ & $1$ & $2$ & $3$ & $4$ & $5$\\\hline
  $P(X=k)$ & & & & & &\\\hline
  \end{tabular}
  \end{center}
  \item Quelle est la valeur de $k$ la plus probable ?
\end{enumerate}

\exo[Des $6$ en quatre lancers]
On lance quatre fois un dé équilibré à six faces et on note $X$ le nombre de $6$
obtenus.
\begin{enumerate}
  \item Justifier que $X$ suit une loi binomiale et préciser ses paramètres.
  \item Calculer $P(X=1)$, sous forme de fraction irréductible.
  \item Calculer $P(1\leq X<3)$, sous forme de fraction irréductible.
\end{enumerate}

\exo[Retrouver un paramètre]
Une variable aléatoire $X$ suit la loi binomiale $\mathcal{B}(3\,;p)$, et
$P(X=0)=\dfrac{8}{125}$.
\begin{enumerate}
  \item Déterminer $p$.
  \item Compléter le tableau de la loi de $X$, sous forme de fractions.
\end{enumerate}

\exo[Des guirlandes de couleur]
Une terrasse est décorée de six guirlandes lumineuses. À l'allumage, chacune prend au
hasard l'une des quatre couleurs rouge, bleu, vert, rose, indépendamment des autres. On
note $X$ le nombre de guirlandes qui s'allument en vert.
\begin{enumerate}
  \item Quelle est la loi de $X$ ? Justifier.
  \item Calculer la probabilité qu'au plus une guirlande s'allume en vert (valeur exacte,
        puis arrondie au millième).
  \item En déduire la probabilité qu'au moins deux guirlandes s'allument en vert.
\end{enumerate}

\exo[D'après bac 2025 et 2026 : vrai ou faux ?]
Pour chaque affirmation, indiquer si elle est vraie ou fausse, en justifiant.
\begin{enumerate}
  \item Dans un musée, la probabilité qu'un visiteur visite sans audioguide est $0{,}32$.
        On choisit au hasard $12$ visiteurs ; le choix de l'audioguide est indépendant
        d'un visiteur à l'autre.

        \emph{Affirmation} : la probabilité qu'exactement la moitié de ces visiteurs opte
        pour l'audioguide est égale à $924\times 0{,}2176^{6}$.
  \item Soit $p$ un nombre réel tel que $0<p<1$, et $X$ une variable aléatoire qui suit
        la loi binomiale de paramètres $3$ et $p$.

        \emph{Affirmation} : $P(X=1)=3p-6p^{2}+3p^{3}$.
\end{enumerate}

\partiefiche{Probabilités cumulées à la calculatrice}

\exo[Traduire, puis calculer]
Soit $X$ une variable aléatoire qui suit la loi binomiale $\mathcal{B}(20\,;0{,}3)$. Pour
chaque probabilité, écrire la traduction avec $P(X\leq\dots)$, puis donner une valeur
arrondie au millième.
\begin{enumerate}[label=\alph*)]
  \item $P(X\leq 5)$
  \item $P(X<5)$
  \item $P(X\geq 8)$
  \item $P(X>8)$
  \item $P(4\leq X\leq 9)$
\end{enumerate}

\exo[Des colis abîmés]
Un service de livraison constate que $3\,\%$ des colis arrivent abîmés. Un livreur
achemine $40$ colis ; on assimile ces colis à un tirage avec remise, et on note $X$ le
nombre de colis abîmés parmi eux.
\begin{enumerate}
  \item Quelle est la loi de $X$ ? Justifier.
  \item Calculer $P(X=0)$, puis la probabilité qu'au moins un colis soit abîmé.
  \item Calculer la probabilité qu'au plus deux colis soient abîmés.
  \item Calculer la probabilité que plus de trois colis soient abîmés.
\end{enumerate}
Arrondir les résultats au millième.

\exo[Un QCM rempli au hasard]
Un QCM comporte $20$ questions ; pour chacune, quatre réponses sont proposées, dont une
seule est juste. Un candidat répond au hasard à chaque question, indépendamment. On note
$X$ le nombre de réponses justes.
\begin{enumerate}
  \item Quelle est la loi de $X$ ? Justifier.
  \item Calculer la probabilité qu'il ait exactement $5$ réponses justes.
  \item Calculer la probabilité qu'il ait au plus $3$ réponses justes.
  \item Calculer la probabilité qu'il ait au moins la moitié des réponses justes.
        Commenter.
\end{enumerate}
Arrondir les résultats au millième.

\exo[D'après bac 2025, Amérique du Sud : vrai ou faux ?]
Un nageur parcourt le $50$ mètres nage libre en moins de $25$ secondes avec une probabilité
de $0{,}85$. Il effectue, sur une journée, $20$ parcours chronométrés ; on note $X$ le
nombre de fois où il nage cette distance en moins de $25$ secondes. On admet que $X$ suit
la loi binomiale de paramètres $n=20$ et $p=0{,}85$.

\emph{Affirmation} : sachant qu'il a atteint au moins $15$ fois son objectif, une valeur
approchée à $10^{-3}$ de la probabilité qu'il l'ait atteint au moins $18$ fois est
$0{,}434$. Vraie ou fausse ? Justifier.

\exo[Des vélos en libre-service]
Une station compte $50$ vélos en libre-service. À $8$~h, chaque vélo est disponible avec
une probabilité de $0{,}7$, indépendamment des autres. On note $X$ le nombre de vélos
disponibles à $8$~h.
\begin{enumerate}
  \item Quelle est la loi de $X$ ?
  \item Calculer la probabilité qu'entre $30$ et $40$ vélos (inclus) soient disponibles.
  \item Calculer la probabilité qu'au moins $40$ vélos soient disponibles.
  \item Calculer la probabilité qu'au plus $30$ vélos soient disponibles.
\end{enumerate}
Arrondir les résultats au millième.

\partiefiche{Problèmes de seuil}

\exo[Des paniers à trois points]
Une joueuse de basket réussit un tir à trois points avec une probabilité de $0{,}2$ ;
ses tirs sont indépendants. Soit $n\in\Ns$ ; on note $X_{n}$ le nombre de tirs réussis sur
$n$ tirs.
\begin{enumerate}
  \item Quelle est la loi de $X_{n}$ ?
  \item Pour $n=5$, calculer la probabilité qu'elle réussisse au moins un tir, arrondie au
        millième.
  \item Exprimer en fonction de $n$ la probabilité $p_{n}$ qu'elle réussisse au moins un tir
        sur $n$.
  \item Déterminer le plus petit nombre de tirs $n$ pour lequel $p_{n}\geq 0{,}95$.
\end{enumerate}

\exo[Des fleurs blanches]
Dans un sachet de semences, $3\,\%$ des graines donnent une fleur blanche. On sème $n$
graines ($n\in\Ns$), en assimilant ce choix à un tirage avec remise.
\begin{enumerate}
  \item Exprimer en fonction de $n$ la probabilité $p_{n}$ d'obtenir au moins une fleur
        blanche.
  \item Combien de graines faut-il semer, au minimum, pour que $p_{n}\geq 0{,}9$ ?
\end{enumerate}

\exo[Des lots de roses remboursés]
Une roseraie livre des fleuristes par lots de $120$ roses ; $5\,\%$ des roses livrées sont
fanées. On choisit un lot au hasard, assimilé à un tirage avec remise, et on note $X$ le
nombre de roses fanées du lot.
\begin{enumerate}
  \item Quelle est la loi de $X$ ?
  \item La roseraie rembourse un lot lorsqu'il contient strictement plus de $8$ roses
        fanées. Calculer la probabilité qu'un lot soit remboursé, arrondie au centième.
  \item La roseraie veut rembourser un lot seulement s'il contient strictement plus de
        $k$ roses fanées, avec $k$ choisi pour que la probabilité de remboursement soit
        au plus $0{,}01$. Justifier que cela revient à chercher le plus petit entier $k$
        tel que $P(X\leq k)\geq 0{,}99$.
  \item Compléter la fonction ci-dessous pour qu'elle renvoie ce plus petit entier $k$,
        puis donner ce nombre.
\end{enumerate}
\begin{python}
from math import comb

def seuil(s):
    k = 0
    P = 0.95**120
    while ............:
        k = k + 1
        P = P + ............
    return k
\end{python}

\exo[Prévoir les repas]
Une association organise un repas solidaire ; $120$ personnes sont inscrites, et chacune
vient avec une probabilité de $0{,}9$, indépendamment des autres. On note $X$ le nombre de
personnes présentes, et on prépare $k$ repas.
\begin{enumerate}
  \item Quelle est la loi de $X$ ?
  \item Traduire par une inégalité sur $P(X\leq k)$ la condition : \og la probabilité de
        manquer de repas est au plus $0{,}02$\fg.
  \item Déterminer le plus petit nombre de repas à préparer.
\end{enumerate}

\partiefiche{Intervalle de probabilité donnée}

\exo[Un intervalle pour $\mathcal{B}(50\,;0{,}4)$]
Soit $X$ une variable aléatoire qui suit la loi binomiale $\mathcal{B}(50\,;0{,}4)$.
\begin{enumerate}
  \item Déterminer le plus petit entier $a$ tel que $P(X\leq a)>0{,}025$.
  \item Déterminer le plus petit entier $b$ tel que $P(X\leq b)\geq 0{,}975$.
  \item En déduire un intervalle $I$ tel que $P(X\in I)\geq 0{,}95$, puis calculer
        $P(X\in I)$, arrondie au millième.
\end{enumerate}

\exo[Le taux de germination annoncé]
Un producteur de semences affirme que $85\,\%$ de ses graines germent. Un jardinier sème
$200$ graines (choix assimilé à un tirage avec remise) ; $155$ germent. On note $X$ le
nombre de graines qui germent sur $200$, et on suppose l'affirmation exacte.
\begin{enumerate}
  \item Quelle est la loi de $X$ ?
  \item Déterminer deux entiers $a$ et $b$ tels que $P(a\leq X\leq b)\geq 0{,}95$, avec la
        méthode du cours.
  \item Interpréter cet intervalle.
  \item Le jardinier doit-il mettre en doute l'affirmation du producteur ?
\end{enumerate}

\exo[Le sport en club]
Une ligue sportive affirme que $30\,\%$ des habitants d'une ville pratiquent un sport en
club. Un sondage interroge $400$ habitants, choisis au hasard (choix assimilé à un tirage
avec remise) ; $125$ pratiquent un sport en club.
\begin{enumerate}
  \item On suppose l'affirmation exacte ; on note $X$ le nombre de pratiquants parmi les
        $400$ habitants interrogés. Quelle est la loi de $X$ ?
  \item Déterminer deux entiers $a$ et $b$ tels que $P(a\leq X\leq b)\geq 0{,}95$, avec la
        méthode du cours.
  \item Le sondage met-il en doute l'affirmation de la ligue ?
  \item Traduire l'intervalle $\intff{a}{b}$ en fréquences.
\end{enumerate}

\partiefichesansnum{Exercices type bac}

\exo[D'après bac 2025, Amérique du Nord]
\emph{Dans cet exercice, les réponses seront arrondies à $10^{-4}$ près.}

Durant la saison hivernale, la circulation d'un virus a entraîné la contamination de
$2\,\%$ de la population d'un pays. Dans ce pays, $90\,\%$ de la population a été
vaccinée contre ce virus. On constate que $62\,\%$ des personnes contaminées avaient été
vaccinées. On interroge au hasard une personne, et on note les événements $C$ : \og la
personne a été contaminée\fg{} et $V$ : \og la personne a été vaccinée\fg.
\begin{enumerate}
  \item À partir de l'énoncé, donner, sans calcul, les probabilités $P(C)$, $P(V)$ et la
        probabilité conditionnelle $P_{C}(V)$.
  \item
  \begin{enumerate}[label=\alph*)]
    \item Calculer $P(C\cap V)$.
    \item En déduire $P(\overline{C}\cap V)$.
  \end{enumerate}
  \item Construire un arbre de probabilités, de premier niveau $C$, $\overline{C}$ et de
        second niveau $V$, $\overline{V}$, et le compléter.
  \item Calculer $P_{V}(C)$ et interpréter le résultat dans le contexte de l'exercice.
  \item Déterminer si les affirmations suivantes sont vraies ou fausses, en justifiant.
  \begin{enumerate}[label=\alph*)]
    \item \og Parmi les personnes non contaminées, il y a dix fois plus de personnes
          vaccinées que de personnes non vaccinées.\fg
    \item \og Plus de $98\,\%$ de la population vaccinée n'a pas été contaminée.\fg
  \end{enumerate}
  \item On s'intéresse à un échantillon de $20$ personnes choisies au hasard dans la
        population, choix assimilé à des tirages successifs avec remise. On note $X$ la
        variable aléatoire égale au nombre de personnes contaminées ; pour une personne
        choisie au hasard, la probabilité d'être contaminée est $p=0{,}02$.
  \begin{enumerate}[label=\alph*)]
    \item Quelle est la loi suivie par $X$ ? Justifier et donner ses paramètres.
    \item Calculer, en rappelant la formule, la probabilité que $4$ personnes exactement
          soient contaminées dans ce groupe de $20$ personnes.
  \end{enumerate}
\end{enumerate}

\exo[D'après bac 2026, Centres étrangers]
Pour tester la conformité de l'acier d'une lame d'escrime, on la plie puis on la redresse
toutes les secondes jusqu'à la rupture ; la lame est conforme si la rupture intervient
après plus de cinq heures de test.

Un équipementier se fournit auprès de trois fabricants : son stock est composé de
$60\,\%$ de lames du fournisseur A, de $12\,\%$ du fournisseur B, le reste venant du
fournisseur C. Une étude qualité a montré que $90\,\%$ des lames de A, $95\,\%$ des lames de
B et $85\,\%$ des lames de C étaient conformes.

\textbf{Partie A.} Un contrôle est déclenché sur une lame vendue par l'équipementier. On
note $A$, $B$, $C$ les événements \og la lame provient du fournisseur A (B, C)\fg{} et
$T$ : \og la lame est conforme\fg.
\begin{enumerate}
  \item Construire l'arbre de probabilité représentant la situation (premier niveau $A$,
        $B$, $C$ ; second niveau $T$, $\overline{T}$).
  \item Déterminer $P(A\cap\overline{T})$ et interpréter ce résultat.
  \item Démontrer que la probabilité que la lame testée soit conforme est de $0{,}892$.
  \item Sachant que la lame testée n'est pas conforme, déterminer la probabilité qu'elle
        provienne du fournisseur B, arrondie au millième.
\end{enumerate}

\textbf{Partie B.} La probabilité qu'une lame ne soit pas conforme est égale à
$0{,}108$. Lors d'une compétition, l'équipementier apporte un échantillon de $75$ lames de
son stock, choix assimilé à un tirage avec remise. On note $X$ le nombre de lames non
conformes dans cet échantillon.
\begin{enumerate}
  \item On admet que $X$ suit une loi binomiale. Préciser ses paramètres.
  \item Calculer la probabilité que $6$ lames exactement soient non conformes, arrondie
        au millième.
  \item L'équipementier affirme que la probabilité qu'il y ait strictement plus de $8$
        lames non conformes dans cet échantillon est inférieure à $50\,\%$. A-t-il
        raison ?
\end{enumerate}

\exo[D'après bac 2025, Polynésie]
\emph{Les probabilités seront, si nécessaire, arrondies à $10^{-3}$ près.}

Une donnée binaire ne peut prendre que deux valeurs : $0$ ou $1$. Elle est transmise
successivement d'une machine à une autre. Chaque machine transmet la donnée reçue soit de
manière fidèle ($1$ reste $1$ et $0$ reste $0$), soit de façon contraire ($1$ devient $0$ et
$0$ devient $1$). La transmission est fidèle dans $90\,\%$ des cas. La première machine
reçoit toujours la valeur $1$.

\textbf{Partie A.} Pour tout $n\in\Ns$, on note $V_{n}$ : \og la $n$-ième machine détient
la valeur $1$\fg.
\begin{enumerate}
  \item
  \begin{enumerate}[label=\alph*)]
    \item Construire l'arbre de probabilité partant de $V_{1}$, avec $V_{2}$ et
          $\overline{V_{2}}$, puis $V_{3}$ et $\overline{V_{3}}$.
    \item Démontrer que $P(V_{3})=0{,}82$ et interpréter ce résultat.
    \item Sachant que la troisième machine a reçu la valeur $1$, calculer la probabilité
          que la deuxième machine ait aussi reçu la valeur $1$.
  \end{enumerate}
  \item Pour tout $n\in\Ns$, on note $p_{n}=P(V_{n})$ ; ainsi $p_{1}=1$.
  \begin{enumerate}[label=\alph*)]
    \item Démontrer que, pour tout $n\in\Ns$, $p_{n+1}=0{,}8p_{n}+0{,}1$.
    \item Démontrer par récurrence que, pour tout $n\in\Ns$,
          $p_{n}=0{,}5\times 0{,}8^{n-1}+0{,}5$.
    \item Calculer la limite de $p_{n}$ quand $n$ tend vers $+\infty$. Interpréter.
  \end{enumerate}
\end{enumerate}

\textbf{Partie B.} La fonction \texttt{simulation} prend en paramètre le nombre $n$ de
transmissions et renvoie la liste des valeurs détenues par les machines successives ;
l'instruction \texttt{rand()} renvoie un nombre aléatoire de $\intfo{0}{1}$.
\begin{python}
def simulation(n):
    donnee = 1
    liste = [donnee]
    for k in range(n):
        if rand() < 0.1:
            donnee = 1 - donnee
        liste.append(donnee)
    return liste
\end{python}
Par exemple, \texttt{simulation(3)} peut renvoyer \texttt{[1, 0, 0, 1]}.
\begin{enumerate}
  \item Déterminer le rôle des instructions des lignes 5 et 6.
  \item Calculer la probabilité que \texttt{simulation(4)} renvoie la liste
        \texttt{[1, 1, 1, 1, 1]}, puis la probabilité que \texttt{simulation(6)} renvoie
        la liste \texttt{[1, 0, 1, 0, 0, 1, 1]}.
\end{enumerate}

\end{multicols}

\end{document}
